\section{Experimental design and pre-registration}
\label{sec:design}

% Figure 4 -- experimental design: the sealed benchmark, the 18-cell run matrix, and the analysis arms.
% Counts from PHASE3/PHASE3_DECISIONS.json, PHASE3/ANALYSIS/COMPLETENESS.json and PHASE2/FREEZE_eval3.json.
\begin{figure}[t]
\centering
\begin{tikzpicture}[font=\footnotesize, x=1mm, y=1mm,
  gcellp/.style={draw=cEBG, fill=cAccentFill, line width=0.4pt, minimum width=7.2mm, minimum height=4.6mm,
                inner sep=0pt, font=\tiny, text=cInk},
  gcells/.style={gcellp, draw=cAll, fill=cWarnFill},
  gcellx/.style={draw=cNeutral, fill=white, line width=0.3pt, minimum width=7.2mm, minimum height=4.6mm,
                inner sep=0pt, font=\tiny, text=cNone}]

% ---------------- benchmark
\node[blkacc, minimum width=36mm, minimum height=17mm, align=center] (bench) at (20,48)
  {\textbf{eval3}\\[1pt]\scriptsize 268 valid tasks\\[-1pt]\scriptsize 135 families\\[-1pt]\scriptsize sealed before any run};
\node[lbl, text=cNone, align=center, anchor=north] at (20,37)
  {\scriptsize written by Qwen2.5-14B\\[-1pt]\scriptsize from briefs only;\\[-1pt]\scriptsize gold fixed before any text};

% ---------------- the 18-cell matrix
\node[lbl, anchor=south, font=\scriptsize\bfseries] at (84,60) {18 cells $\times$ 12 conditions};
\foreach \d [count=\i from 0] in {greedy, {T=0.7 s1}, {T=0.7 s2}, {T=0.7 s3}}
  \node[lbl, rotate=36, anchor=west, font=\tiny] at (66+\i*8,56) {\d};
\foreach \m [count=\j from 0] in {{Qwen2.5-3B}, {Qwen2.5-7B}, {Qwen3-8B}, {Llama-3.1-8B}, {Qwen2.5-14B}, {Llama-3.2-3B}}
  \node[lbl, anchor=east, font=\tiny] at (62,50-\j*5.4) {\m};
\foreach \j in {0,1,2,3}
  \foreach \i in {0,1,2,3}
    \node[gcellp] at (66.5+\i*8,50-\j*5.4) {};
\node[gcells] at (66.5,50-4*5.4) {};
\node[gcells] at (66.5,50-5*5.4) {};
\foreach \j in {4,5}
  \foreach \i in {1,2,3}
    \node[gcellx] at (66.5+\i*8,50-\j*5.4) {--};
\node[lbl, text=cEBG, anchor=west, font=\tiny] at (99,45) {16 primary cells};
\node[lbl, text=cAll, anchor=west, font=\tiny] at (99,21) {2 secondary cells};
\node[lbl, anchor=north, align=center, font=\tiny, text=cNone] at (84,13)
  {3216 rows per cell $=$ 268 tasks $\times$ 12 conditions\\[-1pt]\textbf{57{,}888 runs}; 0 missing, 0 duplicate, 0 harness error};

% ---------------- analysis arms
\node[blkacc, minimum width=42mm, minimum height=11mm, align=left] (prim) at (138,52)
  {\scriptsize\textbf{Confirmatory} (pre-registered)\\[-1pt]\scriptsize 16 primary cells pooled; 51{,}456 runs\\[-1pt]\scriptsize 4 contrasts, Holm, $\alpha=0.05$};
\node[blkplain, minimum width=42mm, minimum height=10mm, align=left] (sens) at (138,38)
  {\scriptsize\textbf{Sensitivity} (pre-registered)\\[-1pt]\scriptsize 3 alternative item sets:\\[-1pt]\scriptsize 161 / 107 / 232 tasks};
\node[blkplain, minimum width=42mm, minimum height=9mm, align=left] (sec) at (138,26)
  {\scriptsize\textbf{Secondary}, never pooled\\[-1pt]\scriptsize 2 greedy-only cells};
\node[blkwarn, minimum width=42mm, minimum height=10mm, align=left] (expl) at (138,14)
  {\scriptsize\textbf{Exploratory}, post-hoc\\[-1pt]\scriptsize 7 ablation contrasts (BH)\\[-1pt]\scriptsize $+$ 2 platform checks};

\draw[flowacc] (bench.east) -- (57,48);
\draw[flowacc] (110,40) -- (prim.west);
\draw[flow] (110,37) -- (sens.west);
\draw[flow] (110,30) -- (sec.west);
\draw[flowback] (110,24) -- (expl.west);
\end{tikzpicture}

\vspace{1mm}
\caption[Experimental design]{%
\textbf{Experimental design.} A sealed benchmark, a fully crossed run matrix, and four analysis arms that are kept
separate throughout the paper. Every one of the six models ran all twelve conditions on all 268 tasks; the four
primary models ran greedy decoding and three $T=0.7$ seeds, and the two secondary models ran greedy only and are
never pooled with the primary ones. The matrix completed with no missing pair, no duplicate and no harness error,
so the pre-registered missing-run rule never had to be applied. Only the first arm supports a confirmatory claim:
the sensitivity analyses were pre-registered but are reported as sensitivity, and everything in the fourth arm is
exploratory, with the two platform checks additionally post-hoc (\Cref{sec:statistics}).\\
\emph{Accessible description:} three panels left to right --- the sealed benchmark, a six-by-four grid of cells in
which the top four rows are full and the bottom two have only their first column filled, and a stack of four
analysis arms.\\
\emph{Source:} \srcfile{PHASE3/ANALYSIS/COMPLETENESS.json}; \srcfile{PHASE2/FREEZE_eval3.json};
\srcfile{PHASE3/PHASE3_DECISIONS.json}.}
\label{fig:design}
\end{figure}


\subsection{What was frozen, and when}

The mechanism was frozen on 14 November 2026: nine arbiter modules, the runner, the analysis script and the
component evaluator, with three frozen dependencies, recorded as a hash manifest. The evaluation set was sealed on
16 November, its gold fixed before any agent saw it. The analysis was pre-registered and frozen on 17 November,
with its own SHA-256 recorded in the seal manifest, \emph{before} the first confirmatory run. Every job then
refused to start unless the mechanism files matched the freeze, the sealed directory passed its checksum, and the
pre-registration hash matched the recorded value; all 20 job starts of the campaign passed all three.

\subsection{Hypotheses}

The pre-registration fixes four success contrasts and one harm test, in this order, with Holm correction over the
four:

\begin{description}[leftmargin=4.4em,itemsep=1pt,topsep=2pt,font=\normalfont\sffamily]
\item[H-primary] \cond{EBG} yields higher task success than \cond{RT} (enforced record-and-temporal labels).
\item[S1] \cond{EBG} yields higher task success than \cond{T} (advisory temporal labels).
\item[S2] \cond{EBG} yields higher task success than \cond{LG} (LLM arbiter under the same gated policy).
\item[S3] \cond{EBG} yields higher task success than \cond{B} (memory without labels).
\item[S4] \cond{EBG}'s harmful-answer rate is \emph{no higher} than \cond{RT}'s.
\end{description}

S4 is a test the mechanism can fail, not one it can pass: its possible verdicts are ``falsified'' and ``not
falsified''. The pre-registration also named three outcomes that would falsify the contribution ---
\cond{EBG} $-$ \cond{RT} $\le 0$, \cond{EBG} $-$ \cond{LG} $\le 0$, and \cond{EBG}'s harmful rate exceeding
\cond{RT}'s --- and required all of them to be reported whatever their sign.

\subsection{Unit, clustering and test}

The unit of analysis is the task. Success per task is the mean over a model's decodings; contrasts are paired
within model and decoding. Families --- (category, attribute) pairs, 135 of them among the valid items --- are the
clustering unit, because tasks within a family share a template and are not independent. The test is a two-sided
sign-flip permutation of family-level summed differences with families pooled across models, so a family's sign is
flipped for all models together; with 135 families exact enumeration is infeasible and the analysis uses 100{,}000
seeded Monte-Carlo draws, whose smallest attainable estimate is $1\times10^{-5}$
\citep{good2005permutation}. Confidence intervals are 95\% percentile bootstraps over families with 10{,}000
seeded resamples \citep{efron1979bootstrap,field2007clustered}. Holm's step-down procedure controls the familywise
error rate over the four pre-registered contrasts at $\alpha=0.05$ \citep{holm1979simple}; every other contrast is
exploratory and reported with Benjamini--Hochberg adjustment \citep{benjamini1995fdr}.

\subsection{Exclusions and missing runs}

There are \textbf{no run-level exclusions}. Unparsable outputs, harness errors and step-limit hits are failures. A
(task, condition) pair with no row, or with a harness error, is scored as a failure and counted separately; a cell
with more than 5\% missing pairs is re-submitted once, with the runner resuming. In this campaign no pair was
missing and no harness error occurred, so these rules never had to be applied.

\subsection{A documented inconsistency in the frozen document}

The pre-registration is internally inconsistent about the primary item set. Its \S5 states that the primary set is
every mechanically valid item; its \S6 and \S8 still contain an earlier rule under which flagged items --- those
with annotator disagreement or low confidence --- would be excluded from the primary analysis. The \S6 and \S8
wording was not updated when \S5 was revised.

We follow \S5 and the recorded primary-set decision: the primary set is all 268 valid items, and the
double-agreement and flagged subsets are reported as sensitivity analyses. We report this as a defect in the frozen
document rather than correct it silently, and \Cref{tab:sensitivity} gives the numbers a reader who prefers the
\S6 reading would need: the four contrasts are $+14.3$, $+12.7$, $+19.7$ and $+14.7$\pp{} on the 161-task
double-agreement subset, all with intervals above zero and Holm-adjusted $p<0.05$. The conclusions do not depend
on which reading is taken.
